% 第 I 讲 8.8--8.13；来源：SGA 4-1.pdf，第65--86页（止于第9节标题前）。
\SourcePage{65}
\SGAsection{8.8}{函子 $J\to\operatorname{Ind}(C)$ 的指标表示}
\textbf{8.8.1.} 设
\[
\varphi:J\longrightarrow\operatorname{Ind}(C),\qquad
\text{即 }\varphi\in\operatorname{Ob}\mathscr{H}\!om(J,\operatorname{Ind}(C))
\]
\SourcePage{66}
为一个\Term{functor}{函子}，其中 $C$ 是一个 $\mathcal U$-\Term{category}{范畴}。回忆（8.5.4），$\varphi$ 的一个\Term{indexed-representation}{指标表示}，按定义是 $\mathscr{H}\!om(J,C)$ 的一个\Term{ind-object}{ind-对象}
\[
\psi:I\longrightarrow\mathscr{H}\!om(J,C),
\]
其中 $I$ 是一个\Term{essentially-small-category}{本质小}的\Term{filtered-category}{滤过范畴}，并给定了它在 $\mathscr{H}\!om(J,\operatorname{Ind}(C))$ 中的\Term{inductive-limit}{归纳极限}到 $\varphi$ 的一个\Term{isomorphism}{同构}。因此，$\varphi$ 有\Term{indexed-representation}{指标表示}，当且仅当它\Term{isomorphism}{同构}于 $\mathscr{H}\!om(J,\operatorname{Ind}(C))$ 中由\Term{full-subcategory}{满子范畴} $\mathscr{H}\!om(J,C)$ 的\Term{object}{对象}构成的一个\Term{essentially-small-category}{本质小}\Term{filtered}{滤过}\Term{inductive-limit}{归纳极限}。如果每个\Term{functor}{函子} $\varphi:J\to\operatorname{Ind}(C)$ 都有\Term{indexed-representation}{指标表示}，便称\Term{category}{范畴} $J$ 对 $C$ 是\Term{admissible-category}{可容许的}。

由于在 $\operatorname{Ind}(C)$ 中小\Term{filtered}{滤过}\Term{inductive-limit}{归纳极限}是\Term{representable-functor}{可表示}的（8.5.1），$\mathscr{H}\!om(J,\operatorname{Ind}(C))$ 中也如此，且它们“逐个变元”计算。因此，包含\Term{functor}{函子}
\[
\mathscr{H}\!om(J,C)\lhook\joinrel\longrightarrow
\mathscr{H}\!om(J,\operatorname{Ind}(C))\tag{8.8.1.1}
\]
\Term{canonical}{典范}地延拓为一个\Term{functor}{函子}
\[
\operatorname{Ind}(\mathscr{H}\!om(J,C))\longrightarrow
\mathscr{H}\!om(J,\operatorname{Ind}(C))\tag{8.8.1.2}
\]
（8.7.2）。这个\Term{functor}{函子}的\Term{essential-image}{本质像}中的\Term{object}{对象}，恰好就是有\Term{indexed-representation}{指标表示}的 $\varphi$；所以，说 $J$ 对 $C$ 是\Term{admissible-category}{可容许的}，也就是说\Term{functor}{函子} (8.8.1.2) 是\Term{essentially-surjective-functor}{本质满射的}。顺便指出：

\begin{SGAresult}{命题}{8.8.2}
若\Term{category}{范畴} $J$ \Term{equivalence}{等价}于一个\Term{finite-category}{有限范畴}，则\Term{canonical}{典范}\Term{functor}{函子} (8.8.1.2) 是\Term{fully-faithful-functor}{全忠实的}；因而 $J$ 对 $C$ 是\Term{admissible-category}{可容许的}，当且仅当它是一个\Term{equivalence-of-categories}{范畴等价}。
\end{SGAresult}
\begin{proof}
由 8.7.5 a)，只须证明：对 $\mathscr{H}\!om(J,C)$ 中每个小\Term{filtered}{滤过}\Term{inductive-system}{归纳系统} $(\varphi_i)_{i\in I}$，以及 $\mathscr{H}\!om(C,J)$ 的每个\Term{object}{对象} $\varphi$，\Term{canonical}{典范}映射
\[
\varinjlim_i\operatorname{Hom}(\varphi,\varphi_i)
\longrightarrow\operatorname{Hom}(\varphi,\varinjlim_i\varphi_i)
\tag{8.8.2.1}
\]
是\Term{bijection}{双射}；这里 $\varinjlim_i\varphi_i$ 表示在 $\mathscr{H}\!om(J,\operatorname{Ind}(C))$ 中取的\Term{inductive-limit}{归纳极限}。若 $\varphi,\psi$ 是两个\Term{functor}{函子} $J\to C$，则有一个对 $\varphi,\psi$ 为\Term{functor}{函子}性的\Term{set}{集合}的\Term{exact-diagram}{正合图}
\[
\operatorname{Hom}(\varphi,\psi)\longrightarrow
\prod_{X\in\operatorname{Ob}J}\operatorname{Hom}(\varphi(X),\psi(X))
\rightrightarrows
\prod_{\substack{f\in\operatorname{Fl}J\\f:X\to Y}}
\operatorname{Hom}(\varphi(X),\psi(Y)).
\]
由于\Term{filtered}{滤过}\Term{inductive-limit}{归纳极限}与平行双\Term{arrow}{箭头}的\Term{equalizer}{等化子}及\Term{finite-product}{有限积}交换，在 $J$ \Term{finite}{有限}时便得到 (8.8.2.1) 的\Term{bijection}{双射}性。$J$ \Term{equivalence}{等价}于\Term{finite-category}{有限范畴}的情形立即归结为上一情形。
\end{proof}

\begin{SGAresult}{命题}{8.8.3}
设 $\varphi:J\to\operatorname{Ind}(C)$ 是一个\Term{functor}{函子}，且 $J$ \Term{equivalence}{等价}于一个小\Term{category}{范畴}。要使 $\varphi$ 有\Term{indexed-representation}{指标表示}，只须满足下面两个条件；当 $J$ \Term{equivalence}{等价}于\Term{finite-category}{有限范畴}时，它们也是必要的：
\SourcePage{67}

 a) \Term{category}{范畴} $\mathscr{H}\!om(J,C)_{/\varphi}$（由 $\mathscr{H}\!om(J,\operatorname{Ind}(C))$ 中以 $\varphi$ 为靶、以 $\mathscr{H}\!om(J,C)$ 的\Term{object}{对象}为源的\Term{arrow}{箭头}构成）是\Term{filtered-category}{滤过的}。

 b) 对 $J$ 的每个\Term{object}{对象} $j$，\Term{functor}{函子}
\[
\psi\longmapsto\psi(j):\mathscr{H}\!om(J,C)_{/\varphi}
\longrightarrow C_{/\varphi(j)}\tag{8.8.3.1}
\]
是\Term{cofinal-functor}{共尾的}，即满足 8.1.3 的条件 F 1)、F 2)。
\end{SGAresult}
\begin{proof}
假定 a)、b) 成立，证明 $\varphi$ 有\Term{indexed-representation}{指标表示}。显然可假设 $J$ 小。令
\[
I=\mathscr{H}\!om(J,C)_{/\varphi},\tag{8.8.3.2}
\]
由假设它是一个\Term{filtered-category}{滤过范畴}。先假设 $I$ \Term{essentially-small-category}{本质小}，于是 $I$ 到 $\mathscr{H}\!om(J,C)$ 的“包含”给出 $\mathscr{H}\!om(J,C)$ 的一个\Term{ind-object}{ind-对象} $i\mapsto\psi_i$。此外，有一个\Term{canonical-homomorphism}{典范同态}
\[
\varinjlim_i\psi_i\longrightarrow\varphi,\tag{8.8.3.3}
\]
逐个变元地由从\Term{functor}{函子} (8.8.3.1) \Term{derivation}{导出}的\Term{homomorphism}{同态}
\[
\varinjlim_I\psi_i(j)\longrightarrow\varphi(j)
=\varinjlim_{C_{/\varphi(j)}}X
\]
给出。后一个\Term{functor}{函子}\Term{cofinal-functor}{共尾}，故 (8.8.3.3) 是\Term{isomorphism}{同构}，从而得到结论。如果不假设 $I$ \Term{essentially-small-category}{本质小}，只须构造一个\Term{essentially-small-category}{本质小}的\Term{full-subcategory}{满子范畴} $I'$，使将 $\varinjlim_I$ 换为 $\varinjlim_{I'}$ 后 (8.8.3.3) 仍为\Term{isomorphism}{同构}。为此，只须由\Term{functor}{函子} (8.8.3.1) 诱导的所有复合\Term{functor}{函子}
\[
I'\longrightarrow I\longrightarrow C_{/\varphi(j)}
\]
均为\Term{cofinal-functor}{共尾函子}。于是由下述引理得到结论；它的证明由判据 8.1.3 b) 立即可得，留给读者。
\end{proof}

\begin{SGAresult}{引理}{8.8.3.4}
设 $I$ 是一个\Term{filtered-category}{滤过的} $\mathcal U$-\Term{category}{范畴}，
\[
f_j:I\longrightarrow I_j\qquad(j\in J)
\]
是从 $I$ 到\Term{essentially-small-category}{本质小}的\Term{filtered-category}{滤过范畴}的一小族\Term{cofinal-functor}{共尾函子}。则 $I$ 有一个既\Term{filtered}{滤过}又\Term{essentially-small-category}{本质小}的\Term{full-subcategory}{满子范畴} $I'$，使各 $f_j$ 诱导的\Term{functor}{函子}均\Term{cofinal}{共尾}。
\end{SGAresult}
\SourcePage{68}
\begin{proof}[命题 8.8.3 的证明（续）]
最后，在假设 $J$ \Term{equivalence}{等价}于\Term{finite-category}{有限范畴}时证明必要性。利用一个\Term{essentially-small-category}{本质小}的\Term{filtered}{滤过}\Term{index-category}{指标范畴} $I$ 给出的 $\varphi$ 的\Term{indexed-representation}{指标表示}定义了一个\Term{functor}{函子} $\psi$：
\[
\begin{tikzcd}[column sep=large,row sep=large]
I\arrow[r,"\psi"]\arrow[dr,dashed,"\psi_j"']&
\mathscr{H}\!om(J,C)_{/\varphi}\arrow[d,dashed]\\
&C_{/\varphi(j)}.
\end{tikzcd}\tag{8.8.3.5}
\]
只须证明 $\psi$ 以及从它\Term{derivation}{导出}的各个\Term{functor}{函子} $\psi_j$ 都\Term{cofinal-functor}{共尾}：由 8.1.3 b)，这的确蕴涵 $\mathscr{H}\!om(J,C)_{/\varphi}$ \Term{filtered}{滤过}，并由定义 8.1.1 得知 (8.8.3.5) 的竖直\Term{arrow}{箭头}也\Term{cofinal}{共尾}。由 8.8.2，可将 $\varphi$ 看成 $\operatorname{Ind}(E)$ 的\Term{object}{对象}，其中 $E=\mathscr{H}\!om(J,C)$；而由以 $I$ 为指标的 $E$ 的\Term{ind-object}{ind-对象} $\varphi$ \Term{derivation}{导出}的\Term{functor}{函子}
\[
\psi:I\longrightarrow E_{/\varphi}
\]
\Term{cofinal}{共尾}，这是一个一般事实，用 8.1.3 的判据 F 1)、F 2) 立即可核实。同样的理由（将 $E,\varphi$ 换成 $C,\varphi(j)$）说明 $\psi_j$ \Term{cofinal}{共尾}，证明完毕。
\end{proof}

\begin{SGAresult}{注记}{8.8.4}
 a) 当 $J$ \Term{equivalence}{等价}于\Term{finite}{有限}\Term{category}{范畴}时，若 $\varphi$ 有\Term{indexed-representation}{指标表示}，我们已经看到 (8.8.3.5) 是一个\Term{cofinal-functor}{共尾函子}，这蕴涵\Term{category}{范畴} $\mathscr{H}\!om(J,C)_{/\varphi}$ 本身\Term{essentially-small-category}{本质小}。

 b) 假定 $C$ 中\Term{finite}{有限}\Term{inductive-limit}{归纳极限}\Term{representable-functor}{可表示}。则 $E=\mathscr{H}\!om(J,C)$ 中也如此，它们逐个变元计算，并且也是 $\mathscr{H}\!om(J,\widehat C)$ 中、从而也是 $\mathscr{H}\!om(J,\operatorname{Ind}(C))$ 中的\Term{inductive-limit}{归纳极限}。立即可知，8.8.3 的条件 a) 此时自动成立，只须考察条件 b)。我不知道，当进一步假设 $J$ 小时，条件 b) 是否自动成立。
\end{SGAresult}
现在给出本小节的主要结果：
\begin{SGAresult}{命题}{8.8.5}
设 $J$ 是一个\Term{equivalence}{等价}于\Term{finite-category}{有限范畴}的\Term{category}{范畴}，并进一步假定 $J$ 是\Term{rigid-category}{刚性的}，即对每个 $j\in\operatorname{Ob}J$，$j$ 的每个\Term{endomorphism}{自同态}都是恒等\Term{morphism}{态射}。则 $J$ 对 $C$ 是\Term{admissible-category}{可容许的}（无论 $\mathcal U$-\Term{category}{范畴} $C$ 为何），即每个\Term{functor}{函子} $J\to\operatorname{Ind}(C)$ 都有\Term{indexed-representation}{指标表示}。
\end{SGAresult}
\begin{proof}
将 $J$ 换成与之\Term{equivalence}{等价}的\Term{category}{范畴}后，可假定 $J$ 是\Term{reduced-category}{约化的}，即 $J$ 中两个\Term{isomorphism}{同构}的\Term{object}{对象}必相等。于是 $J$ 甚至\Term{finite}{有限}。对 $\operatorname{card}\operatorname{Ob}J$ 归纳；此数 $\leq0$ 的情形显然，所以可假定它 $\geq1$。取 $J$ 的一个极大\Term{object}{对象} $j_\circ$，即对每个\Term{arrow}{箭头} $j_\circ\to j$，都存在一个\Term{arrow}{箭头} $j\to j_\circ$；由于 $J$ \Term{rigidity}{刚性}，
\SourcePage{69}
这蕴涵 $j_\circ\to j$ 是\Term{isomorphism}{同构}，又因 $J$ \Term{reduced}{约化}，得到 $j_\circ=j$。令 $J'$ 为从 $J$ 中删去\Term{object}{对象} $j_\circ$ 得到的\Term{full-subcategory}{满子范畴}，$J''$ 为只含\Term{object}{对象} $j_\circ$ 的\Term{full-subcategory}{满子范畴}。命题于是由以下引理推出。
\end{proof}
\begin{SGAresult}{引理}{8.8.5.1}
设 $J$ 是一个\Term{finite}{有限}\Term{category}{范畴}，$J',J''$ 是两个\Term{full-subcategory}{满子范畴}，满足 $\operatorname{Ob}J=\operatorname{Ob}J'\cup\operatorname{Ob}J''$，且对任意 $j'\in\operatorname{Ob}J'$、$j''\in\operatorname{Ob}J''$ 有 $\operatorname{Hom}(j'',j')=\varnothing$。如果 $J',J''$ 对 $C$ 是\Term{admissible-category}{可容许的}，则 $J$ 也如此。
\end{SGAresult}
\begin{proof}
只须证明 8.8.3 的判据 a)、b)，这归结为六项初等核查：对于 $\mathscr{H}\!om(J,C)_{/\varphi}$，核查\Term{filtered-category}{滤过范畴}的最后两个条件；对于\Term{functor}{函子} (8.8.3.1)，分别在 $j\in\operatorname{Ob}J'$ 和 $j\in\operatorname{Ob}J''$ 的情形核查条件 F 1)、F 2)。后面这些条件还蕴涵 $\mathscr{H}\!om(J,C)_{/\varphi}$ \Term{nonempty}{非空}。核查相当繁琐，但并无困难，留给读者。（执笔者曾试图寻找一个优美的证明，以绕开这些令人不快的核查，却未能如愿。）
\end{proof}
\begin{SGAresult}{练习}{8.8.6}
 a) 设 $J,J'$ 是对 $C$ \Term{admissible-category}{可容许的}两个\Term{category}{范畴}，其中一个\Term{finite}{有限}。证明 $J\times J'$ 对 $C$ \Term{admissible}{可容许}。（使用 8.8.2。）

 b) 证明小的\Term{discrete-category}{离散范畴}对每个\Term{category}{范畴} $C$ 都\Term{admissible}{可容许}。

 c) 设\Term{category}{范畴} $J$ 满足：(i) $J$ \Term{rigid-category}{刚性}；(ii) $\operatorname{Ob}J$ 可数；(iii) 对 $J$ 中任意两个\Term{object}{对象} $j,j'$，$\operatorname{Hom}(j,j')$ \Term{finite}{有限}；(iv) 对 $J$ 的每个\Term{object}{对象} $j$，被 $j$ 控制的\Term{object}{对象}，即以 $j$ 为靶的\Term{arrow}{箭头}的源，构成\Term{finite}{有限}\Term{set}{集合}。证明 $J$ 对每个 $\mathcal U$-\Term{category}{范畴} $C$ 都\Term{admissible}{可容许}。（先证明 $J$ 可写成\Term{finite}{有限}\Term{full-subcategory}{满子范畴} $J_n$ 的\Term{filtered}{滤过}并，且凡被 $J_n$ 的某\Term{object}{对象}控制的 $J$ 的\Term{object}{对象}，都属于 $J_n$。再核查 8.8.3 的判据 a)、b)。）
\end{SGAresult}
\begin{SGAresult}{练习}{8.8.7}
在记号上，将\Term{ordered-set}{有序集}与它所定义的\Term{category}{范畴}等同。

 a) 设 $C$ 是\Term{ordered-set}{有序集}。令 $\widetilde C$ 为 $C$ 的所有满足下列条件的\Term{subset}{子集} $A$ 构成的\Term{set}{集合}：$A$ \Term{filtered}{滤过}，且 $x\in A$、$y\leq x$ 蕴涵 $y\in A$。令 $L_\circ:C\to\widetilde C$ 为将每个 $x\in C$ 映为满足 $y\leq x$ 的 $y\in A$ 所构成的\Term{set}{集合} $L_\circ(x)$ 的映射。证明 $L$ 是\Term{injectivity}{单射}，且 $C$ 的序由 $\widetilde C$ 的序诱导。对每个 $A\in\widetilde C$，考虑包含映射 $f(A):A\to C$，它是 $C$ 的一个\Term{ind-object}{ind-对象}；对 $C$ 的每个 \Term{ind-object}{ind-对象} $\varphi:I\to C$，令 $g(\varphi)\in\widetilde C$ 为被某个形如 $\varphi(i)$ 的元素控制的 $C$ 的元素构成的\Term{subset}{子集}。证明这样得到两个互为\Term{quasi-inverse}{拟逆}的\Term{equivalence-of-categories}{等价}
\[
\operatorname{Ind}(C)\mathrel{\mathop{\rightleftarrows}^{f,g}_{\sim}}\widetilde C,
\]
它们将\Term{functor}{函子} $L_\circ$ 变为\Term{canonical}{典范}\Term{functor}{函子} $L:C\to\operatorname{Ind}(C)$。

 b) 假定对每个 $x\in C$，大于等于（分别小于等于）$x$ 的元素所构成的\Term{set}{集合}\Term{finite}{有限}。证明 $C$ 的每个\Term{ind-object}{ind-对象}（分别\Term{pro-object}{pro-对象}）都是\Term{essentially-constant-object}{本质常值的}；
\SourcePage{70}
甚至对每个这样的 $\varphi:I\to C$（分别 $\varphi:I^\circ\to C$），存在 $i_\circ\in\operatorname{Ob}I$，使在 $I_{/i_\circ}$ 上诱导的\Term{functor}{函子}为常值（即把每个\Term{arrow}{箭头}变为\Term{isomorphism}{同构}）。

 c) 取 $C\subset\mathbb N^\circ\times\mathbb N$ 为满足 $i\geq j$ 的自然数对 $(j,i)$ 构成的\Term{subset}{子集}。证明 $\widetilde{\mathbb N}$ \Term{isomorphism}{同构}于给 $\mathbb N$ 加入一个最大元素 $\infty$ 所得的\Term{ordered-set}{有序集}（依 a) 将 $\mathbb N$ 视为 $\widetilde{\mathbb N}$ 的有序\Term{subset}{子集}）。证明 $\widetilde C$ \Term{isomorphism}{同构}于 $\mathbb N^\circ\times\widetilde{\mathbb N}$。考虑\Term{functor}{函子}
\[
\varphi:\mathbb N^\circ\longrightarrow\widetilde C,\qquad
\varphi(j)=(j,\infty).
\]
证明：

 1) $\varphi$ 的\Term{projective-limit}{投射极限}在 $\widetilde C$ 中不\Term{representable-functor}{可表示}，尽管 $C$ 中的\Term{filtered}{滤过}\Term{projective-limit}{投射极限}\Term{representable-functor}{可表示}（且由 b)，是\Term{essentially-constant-object}{本质常值}的）。

 2) 对每个\Term{functor}{函子} $\psi:\mathbb N^\circ\to C$，有 $\operatorname{Hom}(\psi,\varphi)=\varnothing$，因而更不可能将\Term{functor}{函子} $\varphi$ 表成\Term{indexed-representation}{指标形式}。
\end{SGAresult}
\begin{SGAresult}{练习}{8.8.8}
设 $X=\mathbb N\amalg\mathbb N$ 是两个 $\mathbb N$ 的\Term{coproduct}{和集}，$s$ 为 $X$ 的对称变换；对每个 $i\in\mathbb N$，令 $X_i$ 是 $\mathbb N$ 的\Term{subset}{子集} $\Delta_{i+1}\amalg\Delta_i$（其中 $\Delta_i$ 表示由满足 $j\leq i$ 的 $j$ 构成的 $\mathbb N$ 的\Term{subset}{子集}）。令 $C$ 为 $X$ 的\Term{subset}{子集}\Term{category}{范畴}的\Term{subcategory}{子范畴}，其\Term{object}{对象}是各 $X_i$，\Term{arrow}{箭头}是各 $X_i$ 之间由 $X$ 的恒等映射或其对称变换 $s$ 诱导的映射；$\widetilde C$ 以类似方式定义，但还允许\Term{object}{对象} $X$。

 a) 证明 $\operatorname{Ind}(C)$ \Term{equivalence}{等价}于 $\widetilde C$，且\Term{canonical}{典范}\Term{functor}{函子} $C\to\operatorname{Ind}(C)$ \Term{correspondence}{对应}于包含 $C\to\widetilde C$。证明不存在 $C$ 的\Term{object}{对象} $Y$ 及 $Y$ 的\Term{endomorphism}{自同态} $f$ 构成的对 $(Y,f)$，以及从 $(Y,f)$ 到 $(X,s)$ 的 $\operatorname{Ind}(C)$ 中带\Term{endomorphism}{自同态}\Term{object}{对象}的\Term{morphism}{态射}。

 b) 由此推出，若 $J$ 是只有一个\Term{object}{对象}，且除恒等\Term{arrow}{箭头}外仅有一个二阶\Term{arrow}{箭头}的\Term{category}{范畴}，则由 $(X,s)$ 定义的\Term{functor}{函子} $J\to\operatorname{Ind}(C)$ 没有\Term{indexed-representation}{指标表示}。
\end{SGAresult}

\SGAsection{8.9}{$\operatorname{Ind}(C)$ 的正合性质}
\begin{SGAresult}{命题}{8.9.1}
设 $C$ 是一个 $\mathcal U$-\Term{category}{范畴}。

 a) \Term{canonical}{典范}\Term{functor}{函子} $L$（8.2.4.8）和 $c$（8.4.1）
\[
C\xrightarrow{c}\operatorname{Ind}(C)\xrightarrow{L}\widehat C
\]
与\Term{projective-limit}{投射极限}交换。

 b) 假定 $C$ 中\Term{finite}{有限}\Term{inductive-limit}{归纳极限}\Term{representable-functor}{可表示}，且 $C$ \Term{equivalence}{等价}于一个小\Term{category}{范畴}。则 $\operatorname{Ind}(C)$ 中小\Term{projective-limit}{投射极限}\Term{representable-functor}{可表示}。

 c) 若 $C$ 中小\Term{projective-limit}{投射极限}（分别\Term{finite}{有限}\Term{projective-limit}{投射极限}）\Term{representable-functor}{可表示}，则 $\operatorname{Ind}(C)$ 中也如此。设 $J$ 为\Term{finite}{有限}
\SourcePage{71}
\Term{rigid-category}{刚性范畴}，或\Term{discrete-category}{离散范畴}；若 $C$ 中 $J$ 型\Term{projective-limit}{投射极限}\Term{representable-functor}{可表示}，则 $\operatorname{Ind}(C)$ 中同型的\Term{projective-limit}{投射极限}也\Term{representable-functor}{可表示}。

 d) $\operatorname{Ind}(C)$ 中小\Term{filtered}{滤过}\Term{inductive-limit}{归纳极限}（由 8.5.1，它们\Term{representable-functor}{可表示}）是\Term{left-exact-functor}{左正合的}，即与\Term{finite}{有限}\Term{projective-limit}{投射极限}交换。
\end{SGAresult}
\begin{proof}
 a) $L$ 与\Term{projective-limit}{投射极限}交换，来自它的\Term{fully-faithful-functor}{全忠实性}以及 $\widehat C$ 中\Term{projective-limit}{投射极限}“逐个变元”的计算方式；这种计算方式说明：要验证 $\widehat C$ 中的\Term{morphism}{态射}\Term{projective-system}{投射系统} $F\to F_i$（$i\in I$）使 $F$ 成为各 $F_i$ 的\Term{projective-limit}{投射极限}，只须对每个 $X\in\operatorname{Ob}C$，验证\Term{set}{集合}映射的\Term{projective-system}{投射系统} $\operatorname{Hom}(X,F)\to\operatorname{Hom}(X,F_i)$ 满足相应断言。而 $C\subset\operatorname{Ind}(C)$。

$c$ 与\Term{projective-limit}{投射极限}交换，由 $L$ 与之交换，以及复合 $L\circ c:C\to\widehat C$ 与之交换形式地推出。

 b) 考虑到 a)，此断言\Term{equivalence}{等价}于说，\Term{ind-representable-functor}{ind-可表示}\Term{presheaf}{预层}的每个\Term{projective-limit}{投射极限}都 \Term{ind-representable-functor}{ind-可表示}；这由判据 8.3.3 (v) 立即推出，因为\Term{left-exact-functor}{左正合函子}的\Term{projective-limit}{投射极限}仍\Term{left-exact-functor}{左正合}（这又由“\Term{projective-limit}{投射极限}与\Term{projective-limit}{投射极限}交换”（2.5.0）推出）。

 c) 这形式地来自 $\widehat C$ 的相应性质（3.3），以及 $L$ \Term{conservative-functor}{保守}（因为它\Term{fully-faithful-functor}{全忠实}）且与所考虑类型的\Term{limit}{极限}交换（a) 及 8.5.1）。

 d) 众所周知（参见 2.3），\Term{finite}{有限}\Term{projective-limit}{投射极限}\Term{representable-functor}{可表示}，\Term{equivalence}{等价}于下列类型的\Term{projective-limit}{投射极限}\Term{representable-functor}{可表示}：
\[
\text{空图},\quad\ldots,\quad
\begin{tikzcd}[column sep=small,row sep=small]
\bullet\arrow[dr]&&\bullet\arrow[dl]\\
&\bullet&
\end{tikzcd}
\]
（它们\Term{correspondence}{对应}于某些特殊的\Term{finite}{有限}\Term{ordered-set}{有序集}）；而小\Term{projective-limit}{投射极限}\Term{representable-functor}{可表示}，又\Term{equivalence}{等价}于积和\Term{finite}{有限}\Term{projective-limit}{投射极限}\Term{representable-functor}{可表示}。所以 (c) 中的第二个断言蕴涵第一个断言。

先设 $I$ \Term{finite}{有限}。利用将\Term{functor}{函子} $\varphi:J\to\operatorname{Ind}(C)$ 表成\Term{indexed-representation}{指标形式}（8.5.4）的结果 8.8.5，
\[
\phi:J\times I\longrightarrow C,\tag{8.9.1.1}
\]
$\varprojlim\varphi$ 的存在性便是下述更精确、更一般结果的特例。
\end{proof}
\begin{SGAresult}{推论}{8.9.2}
考虑以\Term{indexed-representation}{指标形式} $\Phi$（8.9.1.1）给出的\Term{functor}{函子} $\varphi:J\to\operatorname{Ind}(C)$，其中 $J$ 是一个\Term{finite}{有限}\Term{category}{范畴}。若对每个 $i\in\operatorname{Ob}I$，偏\Term{functor}{函子} $j\mapsto\Phi(j,i)$ 在 $C$ 中有\Term{representable-functor}{可表示}的\Term{projective-limit}{投射极限}（分别\Term{inductive-limit}{归纳极限}），则 $\varphi$ 在 $\operatorname{Ind}(C)$ 中有
\SourcePage{72}
\Term{representable-functor}{可表示}的\Term{projective-limit}{投射极限}（分别\Term{inductive-limit}{归纳极限}），且有\Term{canonical-isomorphism}{典范同构}
\[
\varprojlim\varphi\simeq
\text{“}\varinjlim_i\text{”}\,\varprojlim_j\Phi(j,i)
\tag{8.9.2.1}
\]
（分别
\[
\varinjlim\varphi\simeq
\text{“}\varinjlim_i\text{”}\,\varprojlim_j\Phi(j,i)）,\tag{8.9.2.2}
\]
其中 $\varprojlim_j$（分别 $\varinjlim_j$）在 $C$ 中计算。
\end{SGAresult}
\begin{proof}
第一式只须使用：$\operatorname{Ind}(C)$ 中\Term{filtered}{滤过}\Term{inductive-limit}{归纳极限}与 $\varprojlim_j$ 交换，而 $c$ 与 $\varprojlim_J$ 交换（8.9.1 d)、a)）：
\[
\begin{aligned}
\varprojlim\varphi
&=\varprojlim_j\text{“}\varinjlim_i\text{”}\,\Phi(j,i)\\
&\xleftarrow{\sim}\text{“}\varinjlim_i\text{”}\,
\varprojlim_j^{\operatorname{Ind}(C)}\Phi(j,i)
\simeq\text{“}\varinjlim_i\text{”}\,\varprojlim_j^C\Phi(j,i).
\end{aligned}
\]
第二式类似地证明，只须利用 $I$ 型\Term{inductive-limit}{归纳极限}与 $J$ 型\Term{inductive-limit}{归纳极限}交换（2.5.0），以及下文 8.9.4 a) 中证明的 $c$ 与 $J$ 型\Term{inductive-limit}{归纳极限}交换这一事实。
\end{proof}
\textit{$J$ \Term{discrete}{离散}的情形。} 这一情形包含在以下更精确的断言中。
\begin{SGAresult}{推论}{8.9.3}
设 $I_\alpha$ 是以一个\Term{small-set}{小集合} $A$ 为指标的一族\Term{essentially-small-category}{本质小}的\Term{filtered-category}{滤过范畴}。对每个 $\alpha\in A$，设 $\mathcal X(\alpha)=(X(\alpha)_{i_\alpha})_{i_\alpha\in I_\alpha}$ 是 $C$ 的一个以 $I_\alpha$ 为指标的\Term{ind-object}{ind-对象}。令 $I$ 为各 $I_\alpha$ 的积\Term{category}{范畴}，且假定对每个 $\mathbf i=(i_\alpha)_{\alpha\in A}\in I$，积 $\prod_{\alpha\in A}X(\alpha)_{i_\alpha}$ 在 $C$ 中\Term{representable-functor}{可表示}。则积 $\prod_{\alpha\in A}\mathcal X(\alpha)$ 在 $\operatorname{Ind}(C)$ 中\Term{representable-functor}{可表示}，且它\Term{canonical-isomorphism}{典范同构}于以 $I$ 为指标、由下式给出的 \Term{ind-object}{ind-对象}：
\[
\prod_{\alpha\in A}\text{“}\varinjlim_{i_\alpha\in I_\alpha}\text{”}
X(\alpha)_{i_\alpha}
\simeq\text{“}\varinjlim_{\mathbf i=(i_\alpha)_{\alpha\in A}\in I}\text{”}
\left(\prod_{\alpha\in A}A(\alpha)_{i_\alpha}\right),
\tag{8.9.3.1}
\]
其中右端的积是在 $C$ 中计算的积。（注意，由于各 $I_\alpha$ \Term{filtered}{滤过}且\Term{essentially-small-category}{本质小}，$I$ 也\Term{filtered}{滤过}且\Term{essentially-small-category}{本质小}。）
\end{SGAresult}
\begin{proof}
事实上，众所周知，在 $\widehat C$ 中计算\Term{limit}{极限}时，上式成立（归结到\Term{set}{集合}\Term{category}{范畴}中的情形，立即可见）。结论于是由 $L$ 与所考察的\Term{limit}{极限}交换（8.9.1 a) 及 8.5.1）推出。
\end{proof}
\begin{SGAresult}{注记}{8.9.4}
8.9.2、8.9.3 的证明更一般地表明：若对每个 $i\in\operatorname{Ob}I$，在 $\operatorname{Ind}(C)$ 中计算的 $\varprojlim_j\Phi(j,i)$（分别 $\varinjlim_j\Phi(j,i)$）\Term{representable-functor}{可表示}，则 $\varprojlim_i\varphi$（分别 $\varinjlim_j\varphi$）也\Term{representable-functor}{可表示}。结合 c) 的论证便得：对来自\Term{finite}{有限}或\Term{discrete}{离散}\Term{ordered-set}{有序集}（或更一般地，对 $C$ \Term{admissible-category}{可容许}（8.3.1））的一个给定\Term{category}{范畴} $J$，$\operatorname{Ind}(C)$ 中 $J$ 型\Term{projective-limit}{投射极限}（分别\Term{inductive-limit}{归纳极限}）\Term{representable-functor}{可表示}，只要（且只要）对每个\Term{functor}{函子} $\varphi:J\to C$，$\varphi$ 在 $\operatorname{Ind}(C)$ 中计算的\Term{projective-limit}{投射极限}（分别\Term{inductive-limit}{归纳极限}）
\SourcePage{73}
\Term{representable-functor}{可表示}。在不带“分别”的情形，由 (a)，这还表示 $C$ 上\Term{representable-functor}{可表示}\Term{presheaf}{预层}的每个 $J$ 型\Term{projective-limit}{投射极限}都\Term{ind-representable-functor}{ind-可表示}。
\end{SGAresult}
\begin{SGAresult}{命题}{8.9.5}
设 $C$ 是一个 $\mathcal U$-\Term{category}{范畴}。

 a) \Term{canonical}{典范}\Term{functor}{函子}
\[
c:C\longrightarrow\operatorname{Ind}(C)
\]
是\Term{right-exact-functor}{右正合的}（因而由 8.9.1 a)，是\Term{exact-functor}{正合的}）。

 b) 若 $C$ 中\Term{finite}{有限}\Term{inductive-limit}{归纳极限}（分别\Term{finite}{有限}\Term{coproduct}{和}）\Term{representable-functor}{可表示}，则 $\operatorname{Ind}(C)$ 中小\Term{inductive-limit}{归纳极限}（分别小和）\Term{representable-functor}{可表示}。设 $J$ 为\Term{finite}{有限}或\Term{discrete}{离散}的\Term{preordered-set}{预序集}；若 $C$ 中 $J$ 型\Term{inductive-limit}{归纳极限}\Term{representable-functor}{可表示}，则 $\operatorname{Ind}(C)$ 中也如此。
\end{SGAresult}
\begin{proof}
 a) 假定在 $C$ 中有 $X\simeq\varinjlim_JX_j$，其中 $J$ 是\Term{finite-category}{有限范畴}，$X$ 及各 $X_j$ 均在 $C$ 中。于是，对 $C$ 的每个\Term{ind-object}{ind-对象} $\mathcal Y=(Y_i)_{i\in I}$，有
\[
\begin{aligned}
\operatorname{Hom}(X,\mathcal Y)
&\simeq\varinjlim_i\operatorname{Hom}(X,Y_i)
\simeq\varinjlim_i\varprojlim_j\operatorname{Hom}(X_j,Y_i)\\
&\overset{2.8}{\simeq}\varprojlim_j\varinjlim_i\operatorname{Hom}(X_j,Y_i)
\simeq\varprojlim_j\operatorname{Hom}(X_j,\mathcal Y),
\end{aligned}
\]
比较首尾两项即得所需结论。

 b) 已在 8.8.4 指出，至少对\Term{finite}{有限}\Term{inductive-limit}{归纳极限}而言，这些断言由 a) 及 8.9.2 得出。要证明无限情形中的结论，可以归结到\Term{coproduct}{和}的情形；和可以解释成\Term{finite}{有限}和的\Term{filtered}{滤过}\Term{inductive-limit}{归纳极限}，于是由 8.5.1 得到结论。
\end{proof}
\begin{SGAresult}{注记}{8.9.6}
已经指出，\Term{functor}{函子} $c$ 一般不与\Term{filtered}{滤过}\Term{inductive-limit}{归纳极限}交换，因而也不与无限\Term{coproduct}{和}交换，这与 $L:\operatorname{Ind}(C)\to\widehat C$ 的情形相异。不过，除与\Term{filtered}{滤过}\Term{inductive-limit}{归纳极限}交换（8.5.1）以外，$L$ 几乎从不与任何其他类型的\Term{inductive-limit}{归纳极限}交换（\Term{initial-object}{始对象}、两个\Term{object}{对象}的和、平行双\Term{arrow}{箭头}的\Term{coequalizer}{余等化子}）。例如，若 $C$ 有\Term{initial-object}{始对象} $\varnothing_C$，由 a)，它也是 $\operatorname{Ind}(C)$ 的\Term{initial-object}{始对象}；但 $\varnothing_C$ 从不是 $\widehat C$ 的\Term{initial-object}{始对象}（即恒取 $\varnothing$ 值的\Term{presheaf}{常预层}），因为 $\operatorname{Hom}(\varnothing_C,\varnothing_C)\ne\varnothing$！
\end{SGAresult}
\begin{SGAresult}{命题}{8.9.7}
设 $f:C\to C'$ 是 $\mathcal U$-\Term{category}{范畴}之间的\Term{functor}{函子}，$J$ 为\Term{finite}{有限}\Term{rigid-category}{刚性范畴}（8.8.5）。若 $C$ 中 $J$ 型\Term{inductive-limit}{归纳极限}（分别\Term{projective-limit}{投射极限}）\Term{representable-functor}{可表示}且 $f$ 与之交换，则 $\operatorname{Ind}(f):\operatorname{Ind}(C)\to\operatorname{Ind}(C')$ 也与这一类型的\Term{limit}{极限}交换。
\end{SGAresult}
\SourcePage{74}
\begin{proof}
这由 8.8.5，以及 8.9.2 中对表为\Term{indexed-representation}{指标形式}的\Term{functor}{函子}在 \Term{ind-object}{ind-对象}\Term{category}{范畴}中计算\Term{finite}{有限}\Term{limit}{极限}的办法立即推出。
\end{proof}
\begin{SGAresult}{推论}{8.9.8}
若 $C$ 中\Term{finite}{有限}\Term{inductive-limit}{归纳极限}（分别\Term{projective-limit}{投射极限}）\Term{representable-functor}{可表示}，且 $f$ \Term{right-exact-functor}{右正合}（分别\Term{left-exact-functor}{左正合}），则 $\operatorname{Ind}(f)$ 也如此；在不带“分别”的情形，$\operatorname{Ind}(f)$ 甚至与任意小\Term{inductive-limit}{归纳极限}交换。
\end{SGAresult}
\begin{proof}
第一个断言来自 8.9.7。第二个断言来自第一个断言以及 $\operatorname{Ind}(f)$ 与\Term{filtered}{滤过}\Term{inductive-limit}{归纳极限}交换（8.6.3）这一事实。
\end{proof}
\begin{SGAresult}{练习}{8.9.9}
设 $C$ 是一个 $\mathcal U$-\Term{category}{范畴}。

 a) 假定 $C$ 中\Term{finite}{有限}\Term{coproduct}{和}\Term{representable-functor}{可表示}。证明若 $C$ 中\Term{finite}{有限}和是\Term{disjoint-coproduct}{不交的}（分别\Term{universal-coproduct}{普适的}）（参见第 II 讲 4.5），则 $\operatorname{Ind}(C)$ 中小和也不交（分别\Term{universal}{普适}）。

 b) 假定 $C$ 中每个\Term{morphism}{态射}都\Term{decomposable}{可分解}成一个\Term{epimorphism}{满态射}后接一个\Term{monomorphism}{单态射}（分别为\Term{effective}{有效}\Term{epimorphism}{满态射}后接\Term{monomorphism}{单态射}，分别为\Term{epimorphism}{满态射}后接\Term{effective}{有效}\Term{monomorphism}{单态射}，分别为\Term{effective}{有效}\Term{epimorphism}{满态射}后接\Term{effective}{有效}\Term{monomorphism}{单态射}）。证明 $\operatorname{Ind}(C)$ 满足相同性质。在最后一种“分别”的情形，便可断言 $\operatorname{Ind}(C)$ 的每个\Term{epimorphism}{满态射}都\Term{effective-morphism}{有效}，每个\Term{monomorphism}{单态射}都\Term{effective}{有效}，每个\Term{bimorphism}{双态射}都是\Term{isomorphism}{同构}；证明此时 $\operatorname{Ind}(C)$ 的每个\Term{epimorphism}{满态射}（分别\Term{monomorphism}{单态射}）都可由 $C$ 中\Term{epimorphism}{满态射}（分别\Term{monomorphism}{单态射}）的\Term{inductive-system}{归纳系统}表示。若进一步假定 $C$ 中每个\Term{epimorphism}{满态射}（分别每个\Term{monomorphism}{单态射}）都\Term{universal-morphism}{普适}，则 $\operatorname{Ind}(C)$ 中相同性质成立。

 c) 假定 $C$ \Term{additive-category}{可加}（分别\Term{abelian-category}{阿贝尔}），则 $\operatorname{Ind}(C)$ 也如此。

 d) 设 $X=\text{“}\varinjlim_I\text{”}\,X_i$ 为 $C$ 的一个\Term{ind-object}{ind-对象}，$R\rightrightarrows X$ 为 $X$ 上的\Term{equivalence-relation}{等价关系}。对每个 $i\in\operatorname{Ob}I$，令 $R_i\rightrightarrows X_i$ 为 $R$ 经 $X_i\to X$ 在 $X_i$（视为 $\operatorname{Ind}(C)$ 的\Term{object}{对象}）上诱导的\Term{equivalence-relation}{等价关系}。（假定 $\operatorname{Ind}(C)$ 中的\Term{fiber-product}{纤维积} $R_i=R\times_{X\times X}(X_i\times X_i)$ 可由 $\operatorname{Ind}(C)$ 的\Term{object}{对象}表示；若 $C$ 中\Term{finite}{有限}\Term{projective-limit}{投射极限}\Term{representable-functor}{可表示}，则此假定成立。）证明要使 $R$ \Term{effective-equivalence-relation}{有效}，只须各 $R_i$ \Term{effective}{有效}。

 e) 设 $X$ 是 $C$ 的一个\Term{object}{对象}，$R$ 是 $c(X)$ 上，即视为 $\operatorname{Ind}(C)$ 的\Term{object}{对象}的 $X$ 上的一个\Term{equivalence-relation}{等价关系}。假定 $C$ 中\Term{fiber-product}{纤维积}\Term{representable-functor}{可表示}。要使 $R$ \Term{effective-equivalence-relation}{有效}，必须有 $R=\text{“}\varinjlim_i\text{”}\,R_i$，其中各 $R_i$ 是 $X$（视为 $C$ 的\Term{object}{对象}）上的\Term{equivalence-relation}{等价关系}；若假定 $C$ 中\Term{equivalence-relation}{等价关系}\Term{effective}{有效}，则这一条件也充分。

 f) 假定 $C$ 中每个\Term{morphism}{态射}都\Term{resolution}{分解}成\Term{effective-morphism}{有效}\Term{epimorphism}{满态射}后接\Term{effective}{有效}\Term{monomorphism}{单态射}，且\Term{finite}{有限}\Term{projective-limit}{投射极限}\Term{representable-functor}{可表示}。设 $X$ 是 $C$ 的一个\Term{object}{对象}，$R$ 是 $X\times X$ 的一个\Term{subobject}{子对象}，将它看成 $\operatorname{Ind}(C)$ 的元素。将 $R$ 写成 $\text{“}\varinjlim\text{”}\,Z_i$，其中 $(Z_i)_{i\in I}$ 是 $C$ 中 $X\times X$ 的\Term{subobject}{子对象}的一个\Term{filtered}{滤过}
\SourcePage{75}
递增族（由 b)，可以这样做）。证明要使 $R$ 成为 $X$ 上的\Term{equivalence-relation}{等价关系}，充要条件是对每个 $i\in\operatorname{Ob}I$，存在 $j\in\operatorname{Ob}I$，使之包含 $s(Z_i)$ 和 $Z_i\circ Z_i$，其中 $s$ 是 $X\times X$ 的对称变换。

 g) 假定 $C$ 中\Term{finite}{有限}\Term{projective-limit}{投射极限}及\Term{finite}{有限}\Term{coproduct}{和}\Term{representable-functor}{可表示}，每个\Term{equivalence-relation}{等价关系}\Term{effective}{有效}，且每个\Term{morphism}{态射}都\Term{resolution}{分解}成\Term{effective-morphism}{有效}\Term{epimorphism}{满态射}后接\Term{effective}{有效}\Term{monomorphism}{单态射}。证明 $\operatorname{Ind}(C)$ 中\Term{equivalence-relation}{等价关系}\Term{universal-morphism}{普适}，当且仅当 $C$ 满足下列条件：

\begin{quote}
ST) 对 $C$ 的每个\Term{object}{对象} $X$ 及 $C$ 中 $X\times X$ 的每个\Term{subobject}{子对象} $Z$，若递归地定义 $C$ 中 $X\times X$ 的\Term{subobject}{子对象}序列 $Z_n$（$n\geq0$）为
\[
Z_\circ=Z,\qquad Z_n=\operatorname{Sup}(Z_n,s(Z_n),Z_n\circ Z_n)
\]
（$\operatorname{Sup}$ 在 $X\times X$ 的\Term{subobject}{子对象}集内取，由对 $C$ 的假设，它存在），则序列 $(Z_n)_{n\geq0}$ 是\Term{stationary-sequence}{最终稳定的}。
\end{quote}

 k) 证明 $\operatorname{Ind}((\mathrm{Ens}))$ 中的\Term{equivalence-relation}{等价关系}不一定\Term{effective-equivalence-relation}{有效}。

\textit{*附注。} 关于 $\operatorname{Ind}(C)$ 成为\Term{topos}{拓扑斯}的判据，参见第 VI 讲 8.9.9。*
\end{SGAresult}
\begin{SGAresult}{练习}{8.9.10}
设 $C$ 是一个 $\mathcal U$-\Term{category}{范畴}。令 $\operatorname{Pro}(C)=(\operatorname{Ind}(C^\circ))^\circ$（参见 8.11）。

 a) 假定 $C$ 中\Term{finite}{有限}\Term{coproduct}{和}（分别小和）\Term{representable-functor}{可表示}。证明若这些和\Term{disjoint-coproduct}{不交}（参见第 II 讲 4.5），则 $\operatorname{Pro}(C)$ 满足相同条件。

 b) 设 $J$ 是一个\Term{small-set}{小集合}，$C$ 中以 $J$ 为指标的\Term{coproduct}{和}\Term{representable-functor}{可表示}，因而 $\operatorname{Pro}(C)$ 中也如此。设 $(X(\alpha))_{\alpha\in J}$ 是 $\operatorname{Pro}(C)$ 的一族元素，其中 $X(\alpha)=\text{“}\varprojlim_{I_\alpha}\text{”}\,X_{i_\alpha}$。证明要使各 $X(\alpha)$ 在 $\operatorname{Ind}(C)$ 中的和\Term{universal-coproduct}{普适}，只须对每个族 $\varprojlim_JX_{i_\alpha}$ 也如此，其中 $(i_\alpha)_{\alpha\in J}\in\prod_{\alpha\in J}I_\alpha$。由此推出，若 $C$ 中 $J$ 型和\Term{universal}{普适}，则 $\operatorname{Pro}(C)$ 中也如此的充要条件是：对上述每个满足 $I_\alpha=I$（对所有 $\alpha\in J$）的族 $(X(\alpha))_{\alpha\in J}$，$\operatorname{Pro}(C)$ 中的\Term{canonical-homomorphism}{典范同态}
\[
\text{“}\varprojlim_I\text{”}\,\coprod_\alpha X(\alpha)_i
\longleftarrow
\text{“}\varprojlim_{I^J}\text{”}\,\coprod_{\alpha\in J}X(\alpha)_{i_\alpha}
\]
是\Term{isomorphism}{同构}。由此推出，$\operatorname{Pro}((\mathrm{Ens}))$ 中 $J$ 型和\Term{universal}{普适}，当且仅当 $J$ \Term{finite}{有限}。
\end{SGAresult}
\begin{SGAresult}{练习}{8.9.11}
设 $C$ 是一个 $\mathcal U$-\Term{category}{范畴}。

 a) 设 $(X_i\to X)_{i\in I}$ 为 $C$ 中的一族\Term{morphism}{态射}。证明要使它在 $\operatorname{Ind}(C)$ 中构成\Term{epimorphic-family}{满态射族}，只须存在一个在 $C$ 中构成\Term{epimorphic-family}{满态射族}的\Term{finite}{有限}子族。证明，若假定 $C$ 中每个以 $X_i$ 为靶的\Term{finite-morphism}{有限态射}族都\Term{resolution}{分解}成一个\Term{epimorphic-family}{满态射族}后接一个\Term{effective-morphism}{有效}\Term{monomorphism}{单态射}（10.5），则此条件也必要。（对后一断言，
\SourcePage{76}
令 $P$ 为按包含排序的 $I$ 的\Term{finite}{有限}\Term{subset}{子集}所构成的\Term{set}{集合}，$X'=(X'_J)_{J\in P}$ 为给定族的各\Term{finite}{有限}子族的像构成的\Term{ind-object}{ind-对象}；最后令
\[
X''=X\amalg_{X'}X=\text{“}\varinjlim\text{”}\,X\amalg_{X'_J}X.
\]
证明两个\Term{canonical}{典范}\Term{morphism}{态射} $X\rightrightarrows X''$ 不同，但在各 $X'_J$ 上重合。）

 b) 假定 $C$ 中每个同靶的\Term{finite-morphism}{有限态射}族都\Term{resolution}{分解}成一个\Term{epimorphic-family}{满态射族}后接一个\Term{effective-morphism}{有效}\Term{monomorphism}{单态射}。证明 $\operatorname{Ind}(C)$ 有一个小的\Term{generating-by-epimorphisms}{由满态射生成的子范畴}（7.1），当且仅当 $C$ 有一个小\Term{subcategory}{子范畴} $C'$，使 $C$ 的每个\Term{object}{对象}都为一个源在 $C'$ 中的\Term{finite}{有限}\Term{epimorphic-family}{满态射族}的靶。若假定 $C$ 有一个小\Term{generating-subcategory}{生成子范畴}（7.1），则 $\operatorname{Ind}(C)$ 有小的\Term{generating-by-epimorphisms}{由满态射生成的子范畴}，当且仅当 $C$ \Term{equivalence}{等价}于一个小\Term{category}{范畴}。（对最后这个断言，使用 7.5.2。）

 c) $\operatorname{Ind}((\mathrm{Ens}))$ 没有小\Term{generating-subcategory}{生成子范畴}。
\end{SGAresult}
\begin{SGAresult}{练习}{8.9.12}
设 $C$ 是一个 $\mathcal U$-\Term{category}{范畴}，考虑 $\operatorname{Pro}(C)=\operatorname{Ind}(C^\circ)^\circ$。

 a) 证明若 $\operatorname{Pro}(C)$ 有一个小的\Term{generating-by-epimorphisms}{由满态射生成的子范畴} $P'$（7.1），则 $C$ 有一个小\Term{subcategory}{子范畴} $C'$，它在 $\operatorname{Pro}(C)$ 中由\Term{epimorphism}{满态射}生成。（取 $P'$ 的\Term{object}{对象}的各\Term{set-component}{分量}所构成的\Term{category}{范畴}。）

 b) 证明\Term{category}{范畴} $\operatorname{Pro}((\mathrm{Ens}))$ 没有小的\Term{generating-by-epimorphisms}{由满态射生成的子范畴}。（若 $C'$ 如 a)，取一个\Term{set}{集合} $X$，其\Term{cardinal}{基数}严格大于 $C'$ 的各元素的\Term{cardinal}{基数}，并考虑由 $X$ 中\Term{cardinal}{基数} $<\operatorname{card}(X)$ 的各\Term{subset}{子集}的补集构成的\Term{pro-object}{pro-集合} $\mathcal X$。证明对每个\Term{cardinal}{基数} $<\operatorname{card}(X)$ 的\Term{nonempty}{非空}\Term{set}{集合} $Y$，有 $\operatorname{Hom}(Y,\mathcal X)=\varnothing$，但 $\mathcal X$ 不\Term{isomorphism}{同构}于常 \Term{pro-set}{pro-集合} $\varnothing$。）
\end{SGAresult}

\SGAsection{8.10}{对偶概念：pro-对象、pro-可表示函子}
设 $C$ 是一个 $\mathcal U$-\Term{category}{范畴}。$C$ 的\Term{pro-object}{pro-对象}是指任意一个\Term{functor}{函子}
\[
\varphi:I^\circ\longrightarrow C,\tag{8.10.1}
\]
其中 $I$ 是一个\Term{essentially-small-category}{本质小}的\Term{filtered-category}{滤过范畴}（称为\Term{index-category}{指标范畴}），且如常以 $I^\circ$ 表示\Term{opposite-category}{对偶范畴}。设 $\mathcal V$ 是满足 $\mathcal V\supset\mathcal U$ 的一个\Term{universe}{宇宙}。注意，将每个以 $I\in\mathcal V$ 为指标的 $C^\circ$ 的\Term{ind-object}{ind-对象} $\psi:I\to C^\circ$ \Term{correspondence}{对应}于 \Term{pro-object}{pro-对象} $\varphi=\psi^\circ:I^\circ\to C$，便得到以 $I\in\mathcal V$ 为指标的 $C$ 的 \Term{pro-object}{pro-对象}与以 $I\in\mathcal V$ 为指标的 $C^\circ$ 的 \Term{ind-object}{ind-对象}之间的一一\Term{correspondence}{对应}。根据这一\Term{correspondence}{对应}，通过“搬运结构并反转\Term{arrow}{箭头}”，由 \Term{ind-object}{ind-对象}之间\Term{morphism}{态射}的相应概念（8.2.4.2）、（8.2.5.1）定义 $C$ 的 \Term{pro-object}{pro-对象}之间\Term{morphism}{态射}的概念；从而定义以 $I\in\mathcal V$ 为指标的 $C$ 的 \Term{pro-object}{pro-对象}\Term{category}{范畴}，以及一个\Term{canonical-isomorphism}{典范同构}
\[
\operatorname{Pro}_{\mathcal V}(C,\mathcal U)
\simeq\operatorname{Ind}_{\mathcal V}(C^\circ,\mathcal U)^\circ;
\tag{8.10.2}
\]
当 $\mathcal V=\mathcal U$ 时，简称为 $C$ 的\Term{pro-object}{pro-对象范畴}，记为 $\operatorname{Pro}_{\mathcal U}(C)$，
\SourcePage{77}
或简记为 $\operatorname{Pro}(C)$，所以它由下式定义：
\[
\operatorname{Pro}(C)=\operatorname{Pro}_{\mathcal U}(C,\mathcal U)
\simeq\operatorname{Ind}(C^\circ)^\circ.\tag{8.10.3}
\]
若两个\Term{pro-object}{pro-对象}以\Term{indexed-representation}{指标形式}给出，
\[
\mathcal X=(X_i)_{i\in I},\qquad\mathcal Y=(Y_j)_{j\in J},\tag{8.10.4}
\]
则公式 (8.2.5.1) 在这里成为
\[
\operatorname{Hom}(\mathcal X,\mathcal Y)
\simeq\varprojlim_j\varinjlim_i\operatorname{Hom}(X_i,Y_j).
\tag{8.10.5}
\]
将\Term{functor}{函子} (8.4.1) 应用于 $C^\circ$，再取\Term{opposite-category}{对偶范畴}，便得一个\Term{canonical}{典范}\Term{functor}{函子}（不会混淆时仍以字母 $c$ 表示），借此可将 $C$ 看成 $\operatorname{Pro}(C)$ 的一个\Term{full-subcategory}{满子范畴}：
\[
c:C\longrightarrow\operatorname{Pro}(C).\tag{8.10.6}
\]
正是为了得到这样的\Term{functor}{函子}，而不是 $C\to\operatorname{Pro}(C)^\circ$，我们才在依据 \Term{ind-object}{ind-对象}\Term{morphism}{态射}的相应定义来定义 \Term{pro-object}{pro-对象}\Term{morphism}{态射}时“反转\Term{arrow}{箭头}”。(8.10.6) 的\Term{essential-image}{本质像}中的 \Term{pro-object}{pro-对象}也称为\Term{essentially-constant-object}{本质常值 pro-对象}。

令
\[
\check C=\widehat{C^\circ}
=\mathscr{H}\!om(C,(\mathrm{Ens})),\tag{8.10.7}
\]
则\Term{functor}{函子} (8.2.4.7) 可看成一个\Term{canonical}{典范}的\Term{fully-faithful-functor}{全忠实函子}
\[
L:\operatorname{Pro}_{\mathcal V}(C,\mathcal U)\longrightarrow\check C^\circ,
\tag{8.10.8}
\]
特别地，有
\[
L:\operatorname{Pro}(C)\lhook\joinrel\longrightarrow\check C^\circ.
\tag{8.10.9}
\]
将前面各小节及后面各小节的结果，按需要从\Term{ind-object}{ind-对象}的语言译成\Term{pro-object}{pro-对象}的语言，留给读者。随数学背景不同，两种语言各有用处；此外还有同时出现这两种概念的情形，例如需要考虑 $\operatorname{Pro}(\operatorname{Ind}(C))$ 或 $\operatorname{Ind}(\operatorname{Pro}(C))$ 这样的复杂\Term{category}{范畴}时。这里只再作几项说明，以固定术语和记号，并阐明这种“运算法”。

\textbf{8.10.10.} 若\Term{covariant-functor}{协变函子} $F:C\to(\mathrm{Ens})$ 属于 (8.10.9)（或 8.10.8，两者\Term{equivalence}{等价}）的\Term{essential-image}{本质像}，便称它为\Term{pro-representable-functor}{pro-可表示函子}；所以 $\check C$ 中由这些\Term{functor}{函子}构成的\Term{full-subcategory}{满子范畴}\Term{equivalence}{等价}于 $\operatorname{Pro}(C)^\circ$。通常更愿意在其\Term{opposite-category}{对偶范畴}中工作，即把 (8.10.9) 的\Term{essential-image}{本质像}看成其中的一个\Term{subcategory}{子范畴}，于是它\Term{equivalence}{等价}于 $\operatorname{Pro}(C)$ 本身。与此用法相适应，往往更适宜将\Term{covariant-functor}{协变函子}
\[
F:C\longrightarrow(\mathrm{Ens})
\]
\SourcePage{78}
视为 $\mathscr{H}\!om(C,(\mathrm{Ens}))^\circ=\check C^\circ$ 的\Term{object}{对象}，并相应地将 $\check C^\circ$ 中\Term{homomorphism}{同态}
\[
X\longrightarrow F
\]
的\Term{category}{范畴}，即由 $C$ 的\Term{object}{对象} $X$ 与 $u\in F(X)$ 组成的对 $(X,u)$ 的\Term{category}{范畴}，记为
\[
{}_F\backslash C.\tag{8.10.11}
\]
按这一约定，便得到一个\Term{covariant-functor}{协变函子}（源\Term{functor}{函子}）
\[
{}_F\backslash C\longrightarrow C,\tag{8.10.12}
\]
而 3.4 改写成
\[
F\xrightarrow{\sim}\text{“}\varprojlim_{({}_F\backslash C)^\circ}\text{”}\,X.
\tag{8.10.13}
\]
\textbf{8.10.14.} 将判据 8.3.3 应用于 $C^\circ$，得到 $F$ 的\Term{pro-representable-functor}{pro-可表示性}判据：充要条件是 $({}_F\backslash C)^\circ$ \Term{filtered-category}{滤过}且\Term{essentially-small-category}{本质小}；若 $C$ \Term{equivalence}{等价}于小\Term{category}{范畴}，则后一条件可省去；若进一步假定 $C$ 中\Term{finite}{有限}\Term{projective-limit}{投射极限}\Term{representable-functor}{可表示}，则上述条件\Term{equivalence}{等价}于 $F$ 与它们交换，即 $F$ \Term{left-exact-functor}{左正合}。

\textbf{8.10.15.} 在 $\operatorname{Pro}(C)$ 中，小\Term{filtered}{滤过}\Term{projective-limit}{投射极限}\Term{representable-functor}{可表示}，且\Term{functor}{函子} (8.10.9) 与之交换；换言之，此\Term{functor}{函子}将 $\operatorname{Pro}(C)$ 中的\Term{projective-limit}{投射极限}变成 $\check C=\mathscr{H}\!om(C,(\mathrm{Ens}))$ 中的\Term{inductive-limit}{归纳极限}，如同它与 (8.10.6) 的复合 $C\to\check C^\circ$ 一样；若把它们看成取值于 $\check C$，两者都有同一个令人不便的性质，即为\Term{contravariant-functor}{反变函子}。注意，从 $C$ 到 $(\mathrm{Ens})$ 的\Term{pro-representable-functor}{pro-可表示函子}，是\Term{representable-functor}{可表示函子}的小\Term{filtered}{滤过}\Term{inductive-limit}{归纳极限}，而不是术语可能使人以为的\Term{projective-limit}{投射极限}。

\SGAsection{8.11}{ind-伴随与 pro-伴随}
\textbf{8.11.1.} 设
\[
f:C\longrightarrow C'\tag{8.11.1.1}
\]
是 $\mathcal U$-\Term{category}{范畴}之间的\Term{functor}{函子}，从而得到\Term{functor}{函子} $F\mapsto F\circ f$
\[
f^*:\widehat{C'}\longrightarrow\widehat C.\tag{8.11.1.2}
\]
如果这个\Term{functor}{函子}将\Term{ind-representable-functor}{ind-可表示函子}变成 \Term{ind-representable-functor}{ind-可表示函子}，便说 $f$ 有\Term{ind-adjoint}{ind-伴随}。由于 $f^*$ 与\Term{inductive-limit}{归纳极限}交换，且 $C'$ 中由 \Term{ind-representable-functor}{ind-可表示函子}构成的\Term{full-subcategory}{满子范畴}在小\Term{filtered}{滤过}\Term{inductive-limit}{归纳极限}下\Term{stable}{稳定}，说 $f$ 有 \Term{ind-adjoint}{ind-伴随}，\Term{equivalence}{等价}于说 $f^*$ 将 $C$（视为 $\widehat C$ 的\Term{full-subcategory}{满子范畴}）的\Term{object}{对象}映成 $C$ 上 \Term{ind-representable-functor}{ind-可表示}的\Term{functor}{函子}，即
\[
\begin{gathered}
X\longmapsto\operatorname{Hom}(f(X),Y')\text{ 是 ind-可表示的，}\\
\text{对每个 }Y'\in\operatorname{Ob}C'.
\end{gathered}\tag{8.11.1.3}
\]
\SourcePage{79}
$f$ 有\Term{ind-adjoint}{ind-伴随}这一条件的另一种表述是：存在一个\Term{functor}{函子}
\[
g:\operatorname{Ind}(C')\longrightarrow\operatorname{Ind}(C),\tag{8.11.1.4}
\]
它本质上由 $f^*$ 诱导，即有一个\Term{bifunctor}{双函子}\Term{isomorphism}{同构}
\[
\operatorname{Hom}_{\operatorname{Ind}(C)}(X,g(\mathcal Y'))
\simeq\operatorname{Hom}_{\operatorname{Ind}(C')}(f(X),\mathcal Y'),
\tag{8.11.1.5}
\]
其中 $X\in\operatorname{Ob}C$，$\mathcal Y'\in\operatorname{Ob}\operatorname{Ind}(C')$。事实上，甚至只须有一个\Term{functor}{函子}
\[
g_\circ:C'\longrightarrow\operatorname{Ind}(C)\tag{8.11.1.6}
\]
以及关于 $X\in\operatorname{Ob}C$、$Y'\in\operatorname{Ob}C'$ 的\Term{bifunctor}{双函子}\Term{isomorphism}{同构}
\[
\operatorname{Hom}_{\operatorname{Ind}(C)}(X,g_\circ(Y'))
\simeq\operatorname{Hom}_{C'}(f(X),Y').\tag{8.11.1.7}
\]
显然，\Term{functor}{函子} $g$（分别 $g_\circ$）由 $f$ 确定到唯一\Term{isomorphism}{同构}；反过来，知道这个 $g$（分别 $g_\circ$），也可重建 $f$，相差一个\Term{canonical-isomorphism}{典范同构}。从 (8.11.1.5) 清楚可见，\Term{functor}{函子} $g$ 与\Term{inductive-limit}{归纳极限}交换，并“延拓”\Term{functor}{函子} $g_\circ$；因此这就用 $g_\circ$ 将它确定到唯一\Term{isomorphism}{同构}（8.7.2）。\Term{functor}{函子} $g$，有时也包括它所延拓的\Term{functor}{函子} $g_\circ$，称为 $f$ 的\Term{ind-adjoint}{ind-伴随函子}（这里不必指明“右”，因为另一个\Term{adjunction}{伴随}若存在，将称为 \Term{pro-adjoint}{pro-伴随}，见下文 8.11.5）。

当然，当 $f$ 有\Term{right-adjoint}{右伴随}
\[
f^{\mathrm{ad}}:C'\longrightarrow C
\]
时，它有一个 \Term{ind-adjoint}{ind-伴随} $g$，后者\Term{canonical-isomorphism}{典范同构}于 $f$ 到 \Term{ind-object}{ind-对象}的\Term{canonical-extension}{典范延拓}：
\[
g\simeq\operatorname{Ind}(f^{\mathrm{ad}}).\tag{8.11.1.8}
\]
因此，\Term{ind-adjoint}{ind-伴随}的概念是\Term{right-adjoint}{右伴随}概念的自然推广。

现在考虑\Term{canonical-extension}{典范延拓}
\[
\operatorname{Ind}(f):\operatorname{Ind}(C)\longrightarrow\operatorname{Ind}(C').
\tag{8.11.1.9}
\]
于是有：
\begin{SGAresult}{命题}{8.11.2}
\Term{functor}{函子} $f:C\to C'$ 有\Term{ind-adjoint}{ind-伴随}，当且仅当\Term{functor}{函子} $\operatorname{Ind}(f)$（8.11.1.9）有\Term{right-adjoint}{右伴随}。后者\Term{canonical-isomorphism}{典范同构}于 \Term{ind-adjoint}{ind-伴随} (8.11.1.4)。
\end{SGAresult}
\begin{proof}
充分性以及最后一个断言，由 \Term{ind-adjoint}{ind-伴随}公式 (8.11.1.5) 显然可得。对于必要性，只须注意：在该公式中对各 $X_i$ 取\Term{projective-limit}{投射极限}，便得到关于 \Term{pro-object}{pro-对象} $\mathcal X=(X_i)_{i\in I}$ 与 $\mathcal Y'$ 的\Term{adjunction-isomorphism}{伴随同构}
\[
\operatorname{Hom}_{\operatorname{Ind}(C)}(\mathcal X,g(\mathcal Y'))
\simeq\operatorname{Hom}_{\operatorname{Ind}(C')}
(\operatorname{Ind}(f)(\mathcal X),\mathcal Y').\tag{8.11.2.1}
\]
\SourcePage{80}
证毕。
\end{proof}
\begin{SGAresult}{推论}{8.11.3}
若 $f$ 有\Term{ind-adjoint}{ind-伴随}，则 $\operatorname{Ind}(f)$ 与\Term{inductive-limit}{归纳极限}交换，而 \Term{ind-adjoint}{ind-伴随} $g$ 与\Term{projective-limit}{投射极限}交换。
\end{SGAresult}
\begin{SGAresult}{命题}{8.11.4}
\Term{functor}{函子} $f:C\to C'$ 有\Term{ind-adjoint}{ind-伴随}的必要条件是 $f$ \Term{left-exact-functor}{左正合}；当 $C$ \Term{equivalence}{等价}于一个小\Term{category}{范畴}时，这个条件也充分。
\end{SGAresult}
这由判据 (8.11.1.3) 以及 8.3.1、8.3.3 (iv) 得出。

关于用\Term{accessible-functor}{可达函子}概念表述的另一个判据，参见 8.13.3。

\textbf{8.11.5.} 现在考虑由 $F$ 诱导的\Term{functor}{函子} $F\mapsto F\circ f$
\[
f^{\circ *}:\check{C'}\longrightarrow\check C.\tag{8.11.5.1}
\]
如果这个\Term{functor}{函子}将\Term{pro-representable-functor}{pro-可表示函子}映为 \Term{pro-representable-functor}{pro-可表示函子}，便说 $f$ 有\Term{pro-adjoint}{pro-伴随}，即存在一个\Term{functor}{函子}（称为 $f$ 的 \Term{pro-adjoint}{pro-伴随}\Term{functor}{函子}）
\[
g:\operatorname{Pro}(C)\longrightarrow\operatorname{Pro}(C'),\tag{8.11.5.2}
\]
以及一个\Term{bifunctor}{双函子}\Term{isomorphism}{同构}
\[
\operatorname{Hom}_{\operatorname{Pro}(C)}(g(\mathcal Y'),X)
\simeq\operatorname{Hom}_{\operatorname{Pro}(C')}(\mathcal Y',f(X)).
\tag{8.11.5.3}
\]
当然，说 $f$ 有 \Term{pro-adjoint}{pro-伴随}，就是说 $f^\circ$ 有\Term{ind-adjoint}{ind-伴随}，所以关于 \Term{ind-adjoint}{ind-伴随}的概念和结果都可直接译成关于 \Term{pro-adjoint}{pro-伴随}的陈述。这里只指出，$f$ 有 \Term{pro-adjoint}{pro-伴随}，当且仅当 $\operatorname{Pro}(f):\operatorname{Pro}(C)\to\operatorname{Pro}(C')$ 有\Term{right-adjoint}{右伴随}；此时上述\Term{functor}{函子} $g$ 就是 $\operatorname{Pro}(C)$ 的这样一个\Term{right-adjoint}{右伴随}；还指出，为此 $f$ 必须\Term{left-exact-functor}{左正合}，当 $C$ \Term{equivalence}{等价}于小\Term{category}{范畴}时，这个条件也充分。在此情形，$f$ \Term{exact-functor}{正合}，当且仅当它同时有 \Term{ind-adjoint}{ind-伴随}和 \Term{pro-adjoint}{pro-伴随}。

\begin{SGAresult}{例}{8.11.6}
考虑\Term{functor}{函子}
\[
f:C\longrightarrow(\mathrm{Ens})
\]
的情形。若此\Term{functor}{函子}有\Term{pro-adjoint}{pro-伴随}，则它\Term{pro-representable-functor}{pro-可表示}；反之成立，当且仅当 $\check C$ 中由 \Term{pro-representable-functor}{pro-可表示函子}构成的\Term{full-subcategory}{满子范畴}在小积下\Term{stable}{稳定}；特别地，若 $C$ \Term{equivalence}{等价}于小\Term{category}{范畴}（8.10.14），或 $C$ 中小积\Term{representable-functor}{可表示}（对 $C^\circ$ 应用 8.9.5 b)），便是如此。因而大致说来，对于\Term{functor}{函子} $f:C\to C'$，存在 \Term{pro-adjoint}{pro-伴随}这一概念，是当 $C'=(\mathrm{Ens})$ 时定义的 $f$ 的 \Term{pro-representable-functor}{pro-可表示}性概念的自然推广。
\end{SGAresult}

\SGAsection{8.12}{严格 ind-对象与 pro-对象；一个可表示性判据的应用}
\SourcePage{81}
\textbf{8.12.1.} 设
\[
\mathcal X=(X_i)_{i\in I}
\]
是 $\mathcal U$-\Term{category}{范畴} $C$ 的一个\Term{ind-object}{ind-对象}，$F$ 为它所 ind-表示的\Term{presheaf}{预层}。立即可见，下述两种说法\Term{equivalence}{等价}：\Term{canonical}{典范}\Term{morphism}{态射}
\[
X_i\longrightarrow F=\text{“}\varinjlim_i\text{”}\,X_i
\]
是 $\operatorname{Ind}(C)$ 的\Term{monomorphism}{单态射}（或\Term{equivalence}{等价}地，是 $\widehat C$ 的\Term{monomorphism}{单态射}，即“逐个变元”的\Term{functor}{函子}\Term{monomorphism}{单态射}）；以及，对 $I$ 的每个\Term{arrow}{箭头} $i\to j$，相应的\Term{transition-morphism}{过渡箭头}
\[
X_i\longrightarrow X_j
\]
是\Term{monomorphism}{单态射}。当这些条件成立，且 $I$ 还是一个\Term{ordered-category}{有序范畴}时，称 $\mathcal X$ 为\Term{strict-ind-object}{严格 ind-对象}。注意，这个条件在 \Term{ind-object}{ind-对象}的\Term{isomorphism}{同构}下并不保持；若一个 \Term{ind-object}{ind-对象}\Term{isomorphism}{同构}于某个\Term{strict-ind-object}{严格 ind-对象}，则称它为\Term{essentially-strict-ind-object}{本质严格的}。若一个\Term{presheaf}{预层}可由\Term{strict-ind-object}{严格 ind-对象} ind-表示，则称它\Term{strictly-ind-representable-functor}{严格 ind-可表示}；所以 $F=\text{“}\varinjlim_i\text{”}\,X_i$ \Term{strictly-ind-representable-functor}{严格 ind-可表示}，当且仅当 $\mathcal X=(X_i)_{i\in I}$ 本质严格。

\textbf{8.12.1.1.} 设 $F$ 是 $C$ 上的一个\Term{presheaf}{预层}，考虑 $C_{/F}$ 中由源在 $C$ 中、且为\Term{monomorphism}{单态射}的\Term{arrow}{箭头} $X\to F$ 构成的\Term{full-subcategory}{满子范畴}。称它为 $F$ 的\Term{representable-subfunctor}{可表示子函子范畴}；它是 $F$ 的\Term{representable-subfunctor}{可表示子函子}\Term{ordered-set}{有序集}\Term{correspondence}{对应}的\Term{category}{范畴}，其序从 $F$ 的\Term{subobject}{子对象}集的序诱导。由定义立即得出：
\begin{SGAresult}{命题}{8.12.2}
$C$ 上的\Term{presheaf}{预层} $F$ \Term{strictly-ind-representable-functor}{严格 ind-可表示}，当且仅当 $F$ 的\Term{representable-subfunctor}{可表示子函子}的（有序）\Term{category}{范畴} $I$ \Term{filtered-category}{滤过}且\Term{essentially-small-category}{本质小}，并且
\[
\varinjlim_I X_i\xrightarrow{\sim}F.\tag{8.12.2.1}
\]
当 $C$ 的每个\Term{object}{对象} $X$ 的 $C$ 中\Term{subobject}{子对象}集都是小集时（例如，若 $C$ 有一个小\Term{generating-subcategory}{生成子范畴}（7.4）），这蕴涵\Term{representable-subfunctor}{可表示子函子}\Term{category}{范畴}甚至是小\Term{category}{范畴}。
\end{SGAresult}
\textbf{8.12.2.2.} 若 $F$ \Term{strictly-ind-representable-functor}{严格 ind-可表示}，则有一种优先的办法用\Term{ind-object}{ind-对象} ind-表示它，且该\Term{object}{对象}还是\Term{strict-ind-object}{严格 ind-对象}：取表示 (8.12.2.1)。若\Term{strict-ind-object}{严格 ind-对象} $\mathcal Y$ \Term{isomorphism}{同构}于 (8.12.2.1) 中对某个合适的 $F$ 出现的形如 $(X_i)$ 的 \Term{ind-object}{ind-对象}，便称它\Term{saturated-ind-object}{饱和}。因此，在唯一\Term{isomorphism}{同构}意义下，只有一个严格\Term{saturated-ind-object}{饱和 ind-对象}与给定\Term{strict-ind-object}{严格 ind-对象}\Term{isomorphism}{同构}，即 8.12.2 中取 $F=\varinjlim\mathcal Y$（在 $\widehat C$ 中的\Term{limit}{极限}）所得到的那个\Term{object}{对象}。

\SourcePage{82}
\textbf{8.12.2.3.} 假定已知可在\Term{set}{集合} $\operatorname{Ob}C_{/F}$ 中找到一个小\Term{cofinal-subcategory}{共尾子集}；特别地，当 $F$ \Term{ind-representable-functor}{ind-可表示}时便是如此。于是同一条件在 8.12.2 所考虑的\Term{full-subcategory}{满子范畴} $I$ 中也成立，故在判据 8.12.2 中可省去 $I$ \Term{essentially-small-category}{本质小}的条件。

\textbf{8.12.3.} 设 $F$ 是 $C$ 上的一个\Term{presheaf}{预层}，考虑 $C_{/F}$ 的一个\Term{object}{对象}
\[
u:X\longrightarrow F.
\]
若对 $u$ 的每个\Term{resolution}{分解}
\[
X\xrightarrow{p}X'\xrightarrow{u'}F,\tag{8.12.3.1}
\]
只要 $p$ 是\Term{strict-epimorphism}{严格满态射}（10.2），就有 $p$ 是\Term{isomorphism}{同构}，便称 $u$（或对 $(X,u)$）\Term{minimal-pair}{极小}。将 $u$ 视为 $F(X)$ 的\Term{object}{对象}（1.4），则说 $u$ 极小，就表示：对每个\Term{strict-epimorphism}{严格满态射} $p:X\to X'$，只要 $u\in\operatorname{Im}(F(p):F(X')\to F(X))$，就有 $p$ 是\Term{isomorphism}{同构}。下述引理的 a) 有助于说明这个概念。

\begin{SGAresult}{引理}{8.12.4}
设 $F$ 是 $C$ 上的一个\Term{presheaf}{预层}。

 a) 假定 $F$ 将\Term{coequalizer}{余等化子}变成\Term{equalizer}{等化子}，并考虑一个\Term{morphism}{态射}
\[
u:X\longrightarrow F,
\]
其中 $X\in\operatorname{Ob}C$。当 $C$ 中平行双\Term{arrow}{箭头}的\Term{coequalizer}{余等化子}\Term{representable-functor}{可表示}时，要使 $u$ 为\Term{monomorphism}{单态射}，只须 $u$ \Term{minimal-pair}{极小}；若 $C$ 中\Term{fiber-product}{纤维积}\Term{representable-functor}{可表示}，这个条件也必要。

 b) 假定 $C$ 中\Term{finite}{有限}\Term{inductive-limit}{归纳极限}\Term{representable-functor}{可表示}。要使 $F$ 的\Term{representable-subfunctor}{可表示子函子}的\Term{subcategory}{子范畴}（8.12.1.1）\Term{filtered-category}{滤过}（不一定小）且以 $F$ 为\Term{inductive-limit}{归纳极限}，只须 $F$ \Term{left-exact-functor}{左正合}，并且每个满足 $X\in\operatorname{Ob}C$ 的\Term{morphism}{态射} $u:X\to F$ 都\Term{decomposable}{可分解}为
\[
X\longrightarrow X'\xrightarrow{u'}F,
\]
其中 $u'$ 极小；若 $C$ 中\Term{fiber-product}{纤维积}\Term{representable-functor}{可表示}，这个条件也必要。

 c) 假定 $C$ 有一个小\Term{generating-subcategory}{生成子范畴}（7.1），且 $F$ 的\Term{representable-subfunctor}{可表示子函子}\Term{category}{范畴} $I$ \Term{filtered}{滤过}，则 $I$ 是小\Term{category}{范畴}。
\end{SGAresult}
\begin{proof}
 a) 假定 $u$ 极小，证明 $u$ 是\Term{monomorphism}{单态射}，即对满足 $uv=uv'$ 的每对平行\Term{arrow}{箭头} $v,v':Y\rightrightarrows X$，有 $v=v'$。事实上，若 $X'=\operatorname{Coker}(v,v')$，则 $u$ \Term{resolution}{分解}为 $X\xrightarrow pX'\xrightarrow{u'}F$（因 $F$ 将\Term{coequalizer}{余等化子}变为\Term{equalizer}{等化子}）；按构造 $p$ 是\Term{strict-epimorphism}{严格满态射}，故 $p$ 为\Term{isomorphism}{同构}，即 $v=v'$。再设 $u$ 为\Term{monomorphism}{单态射}，证明它极小。考虑一个\Term{resolution}{分解} (8.12.3.1)；由于 $p$ 是严格
\SourcePage{83}
\Term{epimorphism}{满态射}，它是\Term{canonical}{典范}平行双\Term{arrow}{箭头} $v,v':X\times_{X'}X\rightrightarrows X$ 的\Term{coequalizer}{余等化子}；又因 $(u'p)v=(u'p)v'$，且 $u'p=u$ 是\Term{monomorphism}{单态射}，所以 $v=v'$，从而 $p$ 为\Term{isomorphism}{同构}。

 b) 充分性：由于 $F$ \Term{left-exact-functor}{左正合}，$C_{/F}$ \Term{filtered-category}{滤过}。已知 $F=\varinjlim_{C_{/F}}X$，所以由 8.1.3 c)，只须证明 $C_{/F}$ 中由 $F$ 的\Term{subobject}{子对象}构成的\Term{full-subcategory}{满子范畴} $I$ 在 $C_{/F}$ 中\Term{cofinal-subcategory}{共尾}；而由 a) 的充分性部分，每个 $C_{/F}$ 的\Term{object}{对象}都被一个“极小”\Term{object}{对象}控制这一假设，恰好保证了这件事。必要性：由于\Term{right-exact-functor}{右正合函子}的每个\Term{filtered}{滤过}\Term{inductive-limit}{归纳极限}仍\Term{right-exact-functor}{右正合}，第一个条件显然必要。第二个条件于是来自 a) 的必要性部分。

 c) 知道 $C'_{/F}$ 的\Term{full-subcategory}{满子范畴} $C'_{/X}$，就知道 $F$ 的一个\Term{representable-subfunctor}{可表示子函子} $X\hookrightarrow F$；这里 $C$ 是 $C$ 的一个固定的小\Term{generating-subcategory}{生成子范畴}。（使用 $I$ \Term{filtered-category}{滤过的}假设。）因 $C'_{/F}$ 小，其\Term{full-subcategory}{满子范畴}构成小集，遂得结论。
\end{proof}
\begin{SGAresult}{命题}{8.12.5}
设 $C$ 是一个 $\mathcal U$-\Term{category}{范畴}，其中\Term{finite}{有限}\Term{inductive-limit}{归纳极限}\Term{representable-functor}{可表示}，且有一个小\Term{generating-subcategory}{生成子范畴}（7.1）。设 $F$ 是 $C$ 上的\Term{presheaf}{预层}。要使 $F$ \Term{strictly-ind-representable-functor}{严格 ind-可表示}，只须满足下列两个条件；若 $C$ 中\Term{fiber-product}{纤维积}\Term{representable-functor}{可表示}，这些条件也必要：

 a) $F$ \Term{left-exact-functor}{左正合}。

 b) 每个满足 $X\in\operatorname{Ob}C$、$u\in F(X)$ 的对 $(X,u)$，都在 $C_{/F}$ 中被一个\Term{minimal-pair}{极小对}（8.12.3）控制，即存在一个\Term{minimal-pair}{极小对} $(X',u')$（$X'\in\operatorname{Ob}C$，$u'\in F(X)$）及\Term{morphism}{态射} $f:X\to X'$，使 $u=F(f)(u')$。
\end{SGAresult}
充分性由 8.12.2 及 8.12.4 b)、c) 得到，必要性由 8.3.1 及 8.12.4 a) 得到。
\begin{SGAresult}{注记}{8.12.6}
当 $F$ 将 $C$ 中的\Term{pushout}{纤维和}变为\Term{fiber-product}{纤维积}时，立即可见：对如 b) 中给定的对 $(X,u)$，满足 $u\in\operatorname{Im}(F(X')\to F(X))$ 的 $X$ 的\Term{strict-quotient}{严格商对象} $X'$ 构成的\Term{set}{集合}是向下\Term{filtered-category}{滤过的}。这蕴涵，若 $X$ 的\Term{strict-quotient}{严格商对象}集为\Term{artinian-order}{阿廷的}，则 8.12.5 的条件 b) 自动成立。因此，如果对 $C$ 的每个\Term{object}{对象} $X$ 都满足上述关于 $X$ 的条件（此时有时也说 $C^\circ$ 的\Term{object}{对象}为阿廷的），则由 8.12.5，$F$ \Term{strictly-ind-representable-functor}{严格 ind-可表示}，当且仅当 $F$ \Term{left-exact-functor}{左正合}；此时，只要 $F$ \Term{ind-representable-functor}{ind-可表示}，便\Term{strictly-ind-representable-functor}{严格 ind-可表示}（8.3.1）。
\end{SGAresult}
\begin{SGAresult}{推论}{8.12.7}
设 $C$ 是一个 $\mathcal U$-\Term{category}{范畴}，其中小\Term{projective-limit}{投射极限}\Term{representable-functor}{可表示}，且有一个小\Term{cogenerating-subcategory}{余生成子范畴}（7.9、7.13）。则\Term{functor}{函子} $F:C\to(\mathrm{Ens})$ \Term{representable-functor}{可表示}，当且仅当 $F$ 与小\Term{projective-limit}{投射极限}交换。
\end{SGAresult}
\SourcePage{84}
\begin{proof}
必要性显然。对于充分性，将 8.12.5 应用于\Term{opposite-category}{对偶范畴} $C^\circ$。要先证明 $F$（严格）\Term{pro-representable-functor}{pro-可表示}，只须证明每个满足 $X\in\operatorname{Ob}C$、$u\in F(X)$ 的对 $(X,u)$ 都被一个“极小”对控制。而满足 $u\in\operatorname{Im}(F(X_i)\to F(X))$ 的 $X$ 的\Term{strict-subobject}{严格子对象}族 $(X_i)_{i\in I}$ 是小族（7.5 的\Term{duality}{对偶}形式）。由于 $F$ \Term{left-exact-functor}{左正合}，此族\Term{cofiltered-category}{余滤过}（8.12.6）；又因 $F$ 与小\Term{projective-limit}{投射极限}交换，可见 $X'=\varprojlim_iX_i$ 是此族的最小\Term{object}{对象}。（使用 $F$ \Term{left-exact-functor}{左正合}，故将\Term{monomorphism}{单态射}变为\Term{monomorphism}{单态射}这一事实。）若 $u'\in F(X')$ 是在 $F(X)$ 中的像为 $u$ 的唯一元素，则可见 $(X',u')$ 是控制 $(X,u)$ 的\Term{minimal-pair}{极小对}。这证明 $F$ \Term{pro-representable-functor}{pro-可表示}，结论于是由下述引理推出。
\end{proof}
\begin{SGAresult}{引理}{8.12.8.1}
设 $F:C\to(\mathrm{Ens})$ 为\Term{functor}{函子}，其中 $C$ 是一个小\Term{projective-limit}{投射极限}\Term{representable-functor}{可表示}的 $\mathcal U$-\Term{category}{范畴}。$F$ \Term{representable-functor}{可表示}，当且仅当它\Term{pro-representable-functor}{pro-可表示}且与小\Term{projective-limit}{投射极限}交换。
\end{SGAresult}
\begin{proof}
必要性显然，现证充分性。$F$ \Term{representable-functor}{可表示}，显然就是说 ${}_F\backslash C$ 有一个\Term{initial-object}{始对象}（事实上，${}_F\backslash C$ 的\Term{initial-object}{始对象}恰好是\Term{isomorphism}{同构} $F\to X$，即 $F$ 的表示数据）。由于 $F$ \Term{pro-representable-functor}{pro-可表示}，$({}_F\backslash C^\circ)$ \Term{filtered}{滤过}且\Term{equivalence}{等价}于一个小\Term{category}{范畴}，因此 $X=\varprojlim_{({}_F\backslash C)^\circ}X$ 在 $C$ 中\Term{representable-functor}{可表示}；又因 $F$ 与所考察的\Term{limit}{极限}交换，遂知 $X$ 是 ${}_F\backslash C$ 的\Term{initial-object}{始对象}。证毕。
\end{proof}
\begin{SGAresult}{推论}{8.12.8}
对 $C$ 作 8.12.7 的假设，设 $f:C\to C'$ 是从 $C$ 到一个 $\mathcal U$-\Term{category}{范畴} $C'$ 的\Term{functor}{函子}。$f$ 有\Term{left-adjoint}{左伴随}，当且仅当 $f$ 与小\Term{projective-limit}{投射极限}交换。
\end{SGAresult}
事实上，将 8.12.7 应用于形如 $X\mapsto\operatorname{Hom}(Y',f(X))$ 的复合\Term{functor}{函子}，便显然归结到 8.12.7。

\begin{SGAresult}{例}{8.12.9}
如在 7.13 中指出的，如果 $C$ 是\Term{topological-space}{拓扑空间} $X\in\mathcal U$ 上（或更一般地，一个 $\mathcal U$-\Term{site}{景}上（第 II 讲 (3.0.2)））的\Term{set}{集合}的 $\mathcal U$-\Term{sheaf}{层}\Term{category}{范畴}，则 8.12.7 及 8.12.8 中对 $E$ 的假设成立。下面给出一个富于启发性的例子\footnote{原注 (3)：此例由 H. Bass 提出。}，说明 8.12.7 中存在 $C$ 的小\Term{cogenerating-subcategory}{余生成子范畴} $D$ 的假设并非多余。取 $C$ 为属于 $\mathcal U$ 的群的\Term{category}{范畴}。令 $J$ 为 $C$ 中\Term{simple-group}{单群}的\Term{isomorphism-class}{同构类}所构成的\Term{set}{集合}；对每个 $j\in J$，在类 $j$ 中选取一个\Term{simple-group}{单群} $G_j$。令 $I$ 为 $J$ 的 $\mathcal U$-小\Term{subset}{子集}所构成的\Term{filtered}{滤过}\Term{ordered-set}{有序集}，对 $i\in I$，令 $X_i=\coprod_{j\in i}G_j$。于是各 $X_i$ 在 $E$ 中构成一个\Term{transition-morphism}{过渡态射}为
\SourcePage{85}
\Term{epimorphism}{满态射}的\Term{projective-system}{投射系统} $(X_i)_{i\in I}$，相应的\Term{functor}{函子}
\[
F(X)=\varinjlim_i\operatorname{Hom}(X_i,X)\tag{*}
\]
取值于 $(\mathcal U\text{-}\mathrm{Ens})$（尽管指标集 $I$ 的\Term{cardinal}{基数}显然不属于 $\mathcal U$）。更精确地，证明：对 $C$ 的每个小\Term{full-subcategory}{满子范畴} $C_\circ$，存在 $i_\circ\in I$，使 $F$ 在 $C_\circ$ 上的\Term{restriction}{限制}可由 $X_{i_\circ}$ \Term{representable-functor}{表示}；这就同时证明 $F$ 取值于 $\mathcal U\text{-}\mathrm{Ens}$，以及它与小\Term{projective-limit}{投射极限}交换。要证明此断言，只须注意：对每个 $X\in\operatorname{Ob}C_\circ$，令 $J(X)$ 为满足存在从 $G_j$ 到 $X$ 的非\Term{trivial-bundle}{平凡}\Term{homomorphism}{同态}的 $j\in J$ 所构成的\Term{set}{集合}，则其\Term{cardinal}{基数}必为小\Term{cardinal}{基数}，因为这种\Term{morphism}{态射}必为\Term{monomorphism}{单态射}（$G_j$ 为\Term{simple-group}{单群}）。因此，若 $i_\circ$ 是 $J$ 中所有 $J(X)$（$X\in\operatorname{Ob}C_\circ$）的并，则 $i_\circ$ 是小集，即 $i_\circ\in I$，且它满足要求。另一方面，显然 $F$ 不\Term{representable-functor}{可表示}，因为 $\operatorname{card}I\notin\mathcal U$。由此及 8.12.7，得到属于 $\mathcal U$ 的群的\Term{category}{范畴} $C$ 没有小的满\Term{cogenerating-subcategory}{余生成子范畴}。另一方面，$C$ 的\Term{object}{对象} $\mathbb Z$ 是一个\Term{generator}{生成元}；由 7.12 的证明，于是存在一个由两个元素生成的群 $G$，它不能\Term{embedding}{嵌入}属于 $\mathcal U$ 的群的\Term{category}{范畴} $C$ 的一个\Term{injective-object}{内射对象}中。此外，看来也很可能 $C$ 中除\Term{trivial-bundle}{平凡}群外没有其他\Term{injective-object}{内射对象}。
\end{SGAresult}

\SGAsection{8.13}{pro-可表示函子与可达函子}
\textbf{8.13.1.} 本小节使用下一节中的某些概念和结果，尤其是 9.11 和 9.13，以得到一个\Term{pro-representable-functor}{pro-可表示性}判据；我们将在第 IV 讲 9.16 中顺带用到它。以下用 $C$ 表示满足 9.1 b) 的条件 L 的 $\mathcal U$-\Term{category}{范畴}；例如，若 $C$ 中小\Term{filtered}{滤过}\Term{inductive-limit}{归纳极限}\Term{representable-functor}{可表示}，则该条件成立。
\begin{SGAresult}{命题}{8.13.2}
设 $C$ 如上，$f:C\to(\mathrm{Ens})$ 是一个\Term{functor}{函子}。

 a) 假定 $C$ 的每个\Term{object}{对象}都\Term{accessible-object}{可达}（9.3）。若 $f$ \Term{pro-representable-functor}{pro-可表示}，则 $f$ \Term{accessible-functor}{可达}（9.2）且\Term{left-exact-functor}{左正合}。

 b) 假定 $C$ 中\Term{finite}{有限}\Term{projective-limit}{投射极限}\Term{representable-functor}{可表示}，且 $C$ 有一个\Term{cardinal-filtration}{基数滤过}（9.12）。若 $f$ \Term{accessible}{可达}且\Term{left-exact-functor}{左正合}，则 $f$ \Term{pro-representable-functor}{pro-可表示}。
\end{SGAresult}
\begin{proof}
 a) 对 $C$ 的假设表示从 $C$ 到 $(\mathrm{Ens})$ 的\Term{representable-functor}{可表示}\Term{covariant-functor}{协变函子}\Term{accessible}{可达}。因此，这些\Term{functor}{函子}的每个小\Term{inductive-limit}{归纳极限}也\Term{accessible}{可达}（9.6 (i)），从而每个\Term{pro-representable-functor}{pro-可表示函子}也\Term{accessible}{可达}。
\SourcePage{86}

 b) 由 8.3.3 (iii)，只须证明在 $\operatorname{Ob}({}_F\backslash C)^\circ$ 中存在一个小\Term{cofinal-subcategory}{共尾子范畴}。由假设，存在\Term{cardinal}{基数} $\Pi$，使 $f$ 为 $\Pi$-\Term{accessible-functor}{可达}。设 $(X,u)$，$u\in F(X)$，为 ${}_F\backslash C$ 的一个\Term{object}{对象}。沿用 9.12 c) 的记号，有 $X=\varinjlim_iX_i$，其中 $I$ 为相对于 $\Pi$ \Term{large-relative-to-cardinal}{足够大}的\Term{filtered-category}{滤过范畴}，且各 $X_i$ 属于 $C'=\operatorname{Filt}^{\Pi}(C)$，从而
\[
F(X)\xleftarrow{\sim}\varinjlim_iF(X_i).
\]
这表明小\Term{subcategory}{子范畴} $({}_F\backslash C')^\circ$ 在 $({}_F\backslash C)^\circ$ 中\Term{cofinal}{共尾}，证明完毕。
\end{proof}
\begin{SGAresult}{推论}{8.13.3}
设 $C$ 是一个满足以下条件的 $\mathcal U$-\Term{category}{范畴}：

 a) $C$ 中\Term{finite}{有限}\Term{projective-limit}{投射极限}\Term{representable-functor}{可表示}。

 b) $C$ 中小\Term{filtered}{滤过}\Term{inductive-limit}{归纳极限}\Term{representable-functor}{可表示}。

 c) $C$ 的平行双\Term{arrow}{箭头}\Term{category}{范畴}上的\Term{functor}{函子} $\operatorname{Ker}$ \Term{accessible-functor}{可达}（例如，它与小\Term{filtered}{滤过}\Term{inductive-limit}{归纳极限}交换）。

 d) $C$ 的每个\Term{strict-epimorphism}{严格满态射}都\Term{universally-strict-morphism}{普适严格}（10.2）。

 e) $C$ 存在一个小的\Term{generating-by-strict-epimorphisms}{由严格满态射生成的子范畴}（7.1）。

在这些条件下，\Term{functor}{函子} $f:C\to(\mathrm{Ens})$ \Term{pro-representable-functor}{pro-可表示}，当且仅当它\Term{left-exact-functor}{左正合}且\Term{accessible-functor}{可达}（9.2）。
\end{SGAresult}
事实上，由 9.11 及 9.13，分别可知 8.13.2 a)、b) 中对 $C$ 的条件成立。
\begin{SGAresult}{推论}{8.13.4}
设 $f:C\to C'$ 为满足 8.13.3 的条件 a) 至 e) 的 $\mathcal U$-\Term{category}{范畴}之间的\Term{functor}{函子}。$f$ 有\Term{pro-adjoint}{pro-伴随}，当且仅当 $f$ \Term{left-exact-functor}{左正合}且\Term{accessible-functor}{可达}。
\end{SGAresult}
\begin{proof}
由于\Term{representable-functor}{可表示函子}
\[
h_{Y'}:X'\longmapsto\operatorname{Hom}(Y',X'):
C'\longrightarrow(\mathrm{Ens})
\]
\Term{left-exact-functor}{左正合}且\Term{accessible-functor}{可达}（9.11），若 $f$ 也有这两种性质，则它与上述\Term{functor}{函子}的复合也如此，因而由 8.13.3，这些复合\Term{functor}{函子}\Term{pro-representable-functor}{pro-可表示}，即 $f$ 有\Term{pro-adjoint}{pro-伴随}。反过来，假定 $f$ 有 \Term{pro-adjoint}{pro-伴随}，并设 $C'_1$ 为 $C$ 的一个小\Term{generating-subcategory}{生成子范畴}，$\Pi$ 为使各 $Y'\in\operatorname{Ob}C'$ 都 $\Pi$-\Term{accessible-object}{可达}的\Term{cardinal}{基数}；要证明 $f$ 为 $\Pi$-\Term{accessible}{可达}，只须证明它与各 $h_{Y'}$（$Y'\in\operatorname{Ob}C'$）的复合均如此。而由假设，这些复合 \Term{pro-representable-functor}{pro-可表示}，所以它们\Term{accessible}{可达}（8.13.3）；因而只要取 $\Pi$ \Term{large-relative-to-cardinal}{足够大}，它们就都 $\Pi$-\Term{accessible}{可达}。证毕。
\end{proof}
